\subsection{Deep Image Restoration}
\label{sec:dir}
In general optical image restoration, deep learning-based methods, particularly those based on convolutional neural networks (CNNs), have demonstrated significantly better performance than conventional model-based image restoration methods~\cite{gu2012fast, timofte2013anchored, timofte2015a+, michaeli2013nonparametric, he2010single}. These methods typically learn mapping from low-quality to high-quality images using large-scale paired datasets. Notable CNN-based image restoration networks include the SRCNN~\cite{dongSRCNN2014}, a CNN architecture for image SR, and the DnCNN~\cite{zhangDnCNN2017} for image denoising, along with others~\cite{kim2016accurate, wang2018esrgan, zhang2021plug, zhang2020residual, peng2019dilated, chen2022simple, fu2021model}. Recently, Transformer-based image restoration methods~\cite{liangSwinIRImageRestoration2021, wangUformerGeneralUShaped2022, chenActivatingMorePixels2023} have emerged, following the success of Transformers~\cite{vaswani2017attention} in computer vision. However, these learning methods typically require large-scale training datasets.


\subsection{THz Image Restoration}
Conventional algorithm-based methods for THz image restoration primarily rely on deconvolution procedures~\cite{li2010continuous, ding2010high, xu2014high, wongComputationalImageEnhancement2018, ningResolutionEnhancementTerahertz2019}. These methods depend on the accurate estimation of the beam point spread function (PSF) of THz imaging systems. Owing to the success of deep learning in image restoration, researchers have begun to apply it to THz image restoration has also begun. In particular, researchers have applied the leading methods of deep image restoration to optical images, as introduced in Section~\ref{sec:dir}. Dong et al.~\cite{liTerahertzImageSuperresolution2017a} trained the SRCNN~\cite{dongSRCNN2014} model on 91 natural images. Henceforth, various studies pertaining to for THz image restoration have been conducted ~\cite{leiTerahertzTimedomainSuperresolution2021,wongTrainingAutoencoderbasedOptimizers2019,marinaljubenovicJointDeblurringDenoising2020,hanTerahertzImageRestoration2020,wangTerahertzImageSuperresolution2021,longTerahertzImageSuperresolution2019,zhangRestorationIntegratedCircuit2021,zhangTerahertzImageRestoration2022,luMathematicalDegradationModel2021a,ruanEfficientSubpixelConvolutional2022,maoImageFusionBased2020,suDataDrivenBasedTerahertzImage2023,liAdaptiveTerahertzImage2020,wangComplexZeroshotSuperresolution2022,ljubenovicCNNbasedDeblurringTerahertz2020,yangSuperresolutionReconstructionTerahertz2022,ljubenovicBeamShapeEffectsNoise2022}. 

Lei et al.~\cite{leiTerahertzTimedomainSuperresolution2021} investigated the potential of a graph neural network (GNN) by training it on 8 THz images and testing it on one image.
Wong et al.~\cite{wongTrainingAutoencoderbasedOptimizers2019} investigated the image restoration of frequency modulation continuous wave (FMCW) three-dimensional (3D) images using AutoEncoders, thus providing a novel approach in THz imaging for 3D image restoration.
Marinal et al.~\cite{marinaljubenovicJointDeblurringDenoising2020} investigated a non-deep learning-based approach for THz image restoration that does not require the estimation of the PSF of a beam.
Han et al.~\cite{hanTerahertzImageRestoration2020} were the first to employ ZSSR technique in THz image restoration, which entails  learning directly from an image without requiring additional external data.
Wang et al.~\cite{wangTerahertzImageSuperresolution2021} used a complex CNN (CCNN) to enhance the quality of THz images.
Long et al.~\cite{longTerahertzImageSuperresolution2019} presented a method using a network based on residual connections, which marked a significant advancement in THz image restoration.
Zhang et al.~\cite{zhangRestorationIntegratedCircuit2021} developed a wavelet denoising technique for integrated circuits that reduced noise in THz images.
Zhang et al.~\cite{zhangTerahertzImageRestoration2022} provided a benchmark dataset; however, it was synthetically generated using a PSF and may not be optimized for all THz imaging systems.
Lu et al.~\cite{luMathematicalDegradationModel2021a} developed a mathematical degradation model for restoring THz images.
Ruan et al.~\cite{ruanEfficientSubpixelConvolutional2022} improved the restoration quality of THz images using an efficient sub-pixel convolutional network (ESPCN)~\cite{shi2016real}.
Su et al.~\cite{suDataDrivenBasedTerahertzImage2023} investigated the synthesis of THz images using a data-driven approach.
Li et al.~\cite{liAdaptiveTerahertzImage2020} performed THz image restoration using an adaptive network.
Wang et al.~\cite{wangComplexZeroshotSuperresolution2022} proposed a complex ZSSR (CZSSR) technique that utilizes complex domain data.
Ljubenovic et al.~\cite{ljubenovicCNNbasedDeblurringTerahertz2020} applied CNN-based deblurring to aTHz-TDS system to improve the quality of THz images.
Yang et al.~\cite{yangSuperresolutionReconstructionTerahertz2022} implemented SR to THz images using a residual channel attention network (RCAN)~\cite{zhang2018image} structure.

Recent investigations into THz image restoration have progressed significantly through the application of various deep learning techniques. However, a key challenge is the absence of effective training datasets. Most studies have adopted widely used optical image datasets~\cite{agustsson2017ntire, timofte2017ntire} to emulate the characteristics of actual THz images. In this process, synthetic THz images are created by estimating the PSF of the THz imaging system and applying an isotropic Gaussian blur to the optical images.
However, this approach requires the optimization of the PSF for each THz imaging system. Hence, researchers adopted a random degradation approach that applies various forms of blur. However, the PSF in actual THz imaging systems is not precisely isotropic and varies from system to system. Moreover, THz images have characteristics different from those of conventional optical images. Therefore, models trained on artificially generated data using optical images exhibit performance limitations when directly applied to various THz images.

To address these challenges, we propose a self-supervised learning method using hyperspectral images obtained from a THz-TDS system. This approach does not rely on artificially degrading the image quality for training; instead, it learns the frequency characteristics of the acquired hyperspectral images, thereby overcoming the limitations of conventional training methods. This closely mirrors the actual characteristics of THz imaging systems, thereby enhancing the efficiency and accuracy of THz image restoration.





